% =============================================================================
\section{Terrain Geometry}
\label{sec:terrain}
% =============================================================================

\subsection{Projection with UV retention}
\label{sec:uv-retention}

The height field is a $4096^2$ top-down grid spanning \SI{200}{\meter}. Each
panorama pixel below the horizon is unprojected to a 3D point using calibrated
depth and accumulated into the grid cell containing its $(x,z)$.

The contribution here is what is retained alongside the height. For every cell
that receives a genuine measurement we store the panorama $(u,v)$ it was observed
through, together with a mask distinguishing directly measured cells from those
later filled by interpolation. This costs two float channels and removes an entire
class of error.

The alternative --- re-deriving $(u,v)$ from the final geometry --- is what a
naive implementation does, and it fails in a specific way. The elevation angle of
a ground point is $\varphi = \arcsin(y/d)$, so near the nadir a small error in $y$
produces a large error in $\varphi$ and therefore in panorama row. Since $y$ has
by then passed through reconstruction, erosion and meshing, it is not the value
that was measured. Retaining the observed UV means downstream texturing samples
the correspondence that was actually observed.

\subsection{Constrained harmonic completion}

The projected height field is sparse and uneven: dense near the camera, sparse at
range, and empty wherever an object occluded the ground. Completion is posed as a
constrained boundary-value problem. Cells with direct, high-certainty measurements
become Dirichlet boundary conditions; distant ridge anchors extracted from the sky
silhouette are added as further fixed values; the remainder is solved as a
discrete harmonic problem
\begin{equation}
  \nabla^2 h = 0 \quad \text{on } \Omega_{\text{unknown}},
  \qquad h = h_{\text{obs}} \quad \text{on } \partial\Omega_{\text{obs}},
\end{equation}
using a sparse linear solve. Measured terrain is therefore never modified, and
unobserved terrain is the smoothest surface consistent with it.

\subsection{Ridge anchoring and its assumption}

Terrain beyond depth range has no usable measurement, so its elevation comes from
the sky boundary: for each panorama column, the first non-sky pixel below sky
gives an elevation angle, placed at a fixed radius on the grid. This is a
deliberate scale compromise --- a mountain kilometres away cannot be walkable
terrain in a \SI{200}{\meter} grid, so it is brought close and shrunk, preserving
silhouette shape rather than true distance.

The assumption embedded here is that the sky boundary is the crest of
\emph{ground}. \Cref{sec:results} shows what happens when it is not.

\subsection{De-spoking}

The equirectangular-to-grid projection maps each panorama column onto a radial
line on the grid. Any per-column bias in the depth prediction therefore becomes a
radial streak, and because those streaks lie in the \emph{observed} data --- around
$89\%$ of solve nodes are fixed --- the harmonic solve cannot remove them.

We remove them with a targeted filter that subtracts only the tangential
high-frequency component, computed as a polar round-trip difference
\begin{equation}
  h' = h - \big(P(h) - G_\sigma^{\theta}(P(h))\big),
\end{equation}
where $P$ is the polar resample and $G_\sigma^{\theta}$ an angular Gaussian blur.
Taking the difference cancels the resampling error, which would otherwise bake in
concentric rings: resampling error is around \SI{1.2}{\meter} against a spoke
signal of roughly \SI{0.2}{\meter}. Validated end-to-end, the radial elevation
profile is unchanged (\SI{18.1}{\meter}/\SI{30.7}{\meter}/\SI{21.8}{\meter} at
$r=45/60/\SI{80}{\meter}$ against \SI{18.1}{\meter}/\SI{30.7}{\meter}/%
\SI{21.9}{\meter} before), the cliff mask is bit-identical, and the mean absolute
change is \SI{0.19}{\meter} concentrated in the spoke regions.

\subsection{Erosion refinement}

Harmonic completion is smooth by construction, which is correct as an estimator
and wrong as terrain. Real landforms carry drainage structure. We therefore run
fluvial incision over the completed field using
Landlab~\cite{barnhart2020landlab,hobley2017creative}, restricted to unobserved
cells so measured ground is preserved. Erosion strength is modulated by a
jaggedness map derived from the ridge silhouette, so the invented relief matches
the character of the visible terrain.

\subsection{Meshing}

The height field is triangulated by Poisson-disc sampling with a near-camera
density bias, then Delaunay triangulation, giving natural level-of-detail without
axis-aligned artefacts. Water regions are extracted as a separate surface, and the
terrain beneath is depressed so an animated water plane never clips through its
bed. Disconnected landmasses above a plausibility threshold are extracted as
separate formation meshes. A coarser collision mesh is generated by the same
procedure with sparser parameters.

Each vertex carries a second UV channel holding the retained panorama coordinate
from \cref{sec:uv-retention}. Because that channel is per-vertex and interpolated
across triangles, triangles spanning an occlusion boundary --- where grid-adjacent
vertices observed very different panorama rows --- smear a wide band of image
across a small piece of ground. We detect these geometrically, comparing each
triangle's UV span against the span its own size and distance predict, and collapse
the offending corner. Detection must be geometric rather than a fixed pixel budget,
because a near triangle legitimately spans tens of degrees while the same triangle
at range spans under two.
